%!TEX root = ../LLMSeminar.tex
% ──────────────────────────────────────────────────────────────
%  Section 0 — Why This Talk?
% ──────────────────────────────────────────────────────────────

\section{Why this talk?}\label{sec:why}

\subsection{The landscape has changed}

It would be hard to ignore the fact that, in the past couple of years, AI
has entered our lives in one way or another---whether we want it or not.
Some of you are engaging proactively with these tools; others perhaps
reluctantly.  The purpose of this talk is not to evangelise.  It is to
take a sober look, as research scientists, at one question: \emph{are AI
tools relevant for us, and should we take them seriously?}

I spent the last quarter of 2025 going deep on this question.  My
summary, and what I want to share with you today, is that I think we had
better be engaging with these tools.  How you engage is up to you.  I am
just going to share my journey and the things I found important to be
aware of as we enter this new AI era.  They are not going to go away.

\subsection{GPT-5.2 derives new results in theoretical physics}

On 13 February 2026, a collaboration including researchers at the IAS,
Harvard, Cambridge, Vanderbilt, and OpenAI published a paper
(arXiv:2602.12176) claiming that GPT-5.2 had derived new results in
theoretical physics.  Should we worry that a language model is now doing
the job of high-energy theorists?

The authors are not unknown---they are very serious scientists who
successfully carried out a research result assisted by AI\@.  I am no fan
of OpenAI and their press-release culture; the result is not as amazing
as the press release suggests.  But the paper alone demonstrates that
some of the world's greatest researchers are already engaged heavily with
AI tools, especially in the US, and they find them helpful.

\subsection{Three moments}

I see three key moments in the development of LLMs that matter for
research scientists:

\begin{enumerate}[(1)]
  \item \textbf{November 2022: ChatGPT launches.}  The utility of
    ChatGPT was dramatically more than the previous GPTs.  Everyone
    experimented with it.

  \item \textbf{Q1 2025: Claude Code, Cursor, Copilot.}  Something
    deeply significant happened: coding agents were released.  Claude
    Code launched in February 2025 and was a definite inflection point in
    the development of AI tools.

  \item \textbf{Q4 2025: Opus 4.5, GPT-5.2 cross a threshold.}  The
    next-generation frontier models were significantly better.  Even if
    they never improved beyond this point, we would already have entered
    a new era of utility.  The reality is they are getting better still.
\end{enumerate}

\subsection{Quick survey}

\begin{audience}[Audience poll]
  Raise your hand if you have used:
  \begin{itemize}
    \item ChatGPT, Claude, or Gemini?
    \item An LLM in your IDE to write code?
    \item Coding agents?
  \end{itemize}
\end{audience}

\subsection{A skeptic's confession}

If you tried chat in Q1 2025, you probably dismissed LLMs for research.
I did.  I was a skeptic.  I experimented for research, and I am still
annoyed at the models of Q1 2025.

But here is the central tension:

\begin{keyresult}[The paradox]
  LLMs are \textbf{incredibly capable} but \textbf{maximally
  untrustworthy}.
\end{keyresult}

The evidence is damning:
\begin{itemize}
  \item They produce seductively correct equations that are nonsense.
  \item Push back on a valid proof and the LLM finds a ``bug'' that is
    not there.
  \item They lie, cheat, steal, and flatter.
\end{itemize}
All terrible qualities for reproducible science.

\subsection{Mental model}

\begin{intuition}[A good mental model]
  Think of an LLM as a \emph{slightly unhinged, hypermotivated MSc
  student}.  Extremely productive, occasionally brilliant, but you
  would never trust them unsupervised with anything that matters.
\end{intuition}

\subsection{The gap}

The gap between ``I tried ChatGPT and it hallucinated'' and ``this is a
powerful force multiplier'' is \emph{not} closed by better models.  It is
closed by \textbf{technique} and \textbf{domain expertise}.

What it takes:
\begin{enumerate}[(1)]
  \item Deliberate practice---a different workflow.
  \item A subscription---free chat is not enough.
  \item A coding agent---Claude Code, Cline, or similar.
\end{enumerate}

\subsection{Auto-formalisation and research acceleration}

Auto-formalisation is real.  Writing new science, with correct tooling
and expert guidance, is now possible.  The human provides the ground
truth.

LLMs can be understood.  Let me rebuild the stack from first principles.
